\section{Full Proofs}\label{app:proofs}

\subsection{Mergeable Summaries Background}\label{app:mergeable-background}

The main text introduces mergeable sketches in one paragraph
(Section~\ref{sec:mergeable-sketches}). This subsection gives the
self-contained version of that story. The important pattern is simple:
build a state on each shard, merge states up the tree, and apply the
query/readout only at the root. C-TreePO uses the same pattern, except
that the state can be learned and the local laws audit whether it still
contains the information needed by the task.

\paragraph{Five older ideas.}
Database systems already distinguish aggregates whose fixed-size
intermediate state can be combined before finalization
\citep{GrayEtAl1997DataCube}. Streaming theory formalizes this as
mergeable summaries \citep{Agarwal2013MergeableTODS}; distributed
streaming work studies the symmetric unordered case
\citep{FeldmanEtAl2008MUD}; and functional programming gives the
ordered analogue through list, or free-monoid, homomorphisms
\citep{Gibbons1996ThirdHomomorphism}. Finally, the
sufficient-statistic tradition explains why a compressed state is
useful only relative to a downstream target
\citep{Fisher1922,Blackwell1953}. C-TreePO takes this lineage as its
algebraic target: learn a bounded state, merge it locally, and audit
whether the task readout is preserved.

\paragraph{Mergeable summaries in data systems.}
A \emph{mergeable summary} is a compact state with a merge rule. Given
two shards, one builds a state for each shard, merges those two states,
and obtains a state that is valid for the union of the shards. The
query is then read from that merged state. \citet{Agarwal2013MergeableTODS}
formalize this for a query $Q$ (cardinality, frequency, quantile,
membership, \dots): the state obtained from $u\concat v$ and the state
obtained by merging the states for $u$ and $v$ must answer
$Q(u\concat v)$ equally well, while staying within the stated size
bound. The appeal is operational: a MapReduce, Dataflow, or Spark job
can compute a summary per shard, merge summaries pairwise up a tree,
and read the query off the root without materializing the full stream.
The same state-composition pattern appears in massive unordered
distributed aggregation and in database ``algebraic'' aggregates,
where fixed-size intermediate state is combined before finalization
\citep{FeldmanEtAl2008MUD,GrayEtAl1997DataCube}.

The mergeable-summary literature is deep. HyperLogLog
\citep{FlajoletEtAl2007,HeuleEtAl2013} estimates distinct counts in a
few kilobytes regardless of stream length; Appendix~\ref{app:hll-primer}
gives a self-contained primer. Count-Min
\citep{CormodeMuthukrishnan2005} answers frequency queries under a
space budget. Bloom filters, Theta/KMV sketches, t-digest, KLL, and GK
quantile sketches are part of the same systems lane, with different
exact or approximate merge guarantees. The common thread is that the
state, merge, and readout are engineered before training, so
correctness is usually proved from sketch-specific algebra rather than
inferred from runtime audits. Current Lean-backed instances include
HLL's max-register merge laws, Count-Min's additive state merge,
Misra--Gries/GK-style theorem surfaces, typed
interval/range targets from the mergeable-summary literature, adjacent
$\varepsilon$-kernel geometry when separately cited, and the paper's bigram
and Markov-count sketches.
Separately, the Python benchmark suite wraps Apache DataSketches
implementations for HLL, CPC, Theta, Count-Min, Frequent Items, KLL,
classic quantiles, REQ, t-digest, Tuple, and VarOpt as official
empirical references. CPC, Theta/KMV, Frequent Items, classic
quantiles, REQ, t-digest, Tuple, and VarOpt should therefore be read
as empirical-only rows unless a matching local-law artifact is added.
HLL is the empirical anchor in Appendix~\ref{sec:hll-parity}, not the
definition of the classical class.

\paragraph{The algebraic view.}
Abstractly, a mergeable sketch for a query family indexed by $q \in
\mathcal{Q}$ is a state-level triple $(\mathrm{encode}, \oplus,
\mathrm{query})$ with
\[
\mathrm{encode} : \Strings \to \mathcal{S}, \quad
\oplus : \mathcal{S} \times \mathcal{S} \to \mathcal{S}, \quad
\mathrm{query} : \mathcal{S} \times \mathcal{Q} \to \mathcal{R},
\]
and correctness
\[
\mathrm{query}\big(\mathrm{encode}(u \concat v),\, q\big)
\;\approx\;
\mathrm{query}\big(\mathrm{encode}(u) \oplus \mathrm{encode}(v),\, q\big)
\]
for every pair of shards $u, v$ and every query $q$. The merge
composes \emph{states}, not scalar query answers. Distinct count is
the canonical warning: $|A \cup B|$ is not determined by $|A|$ and
$|B|$ alone, while an HLL register vector can still be merged because
it carries overlap-relevant state. Associativity of $\oplus$ is the
schedule-invariance condition: any binary aggregation tree reaches the
same state, or at least the same readout. Commutativity is specific to
unordered or symmetric inputs. Idempotence is specific to set- or
semilattice-like unions such as HLL or Bloom-filter-style membership
sketches; it is not a generic property of frequency or quantile
summaries.

For ordered text, the corresponding algebraic ideal is a list, or
free-monoid, homomorphism: $h(u \concat v) = h(u) \odot h(v)$ for an
associative operator $\odot$ on summary states
\citep{Gibbons1996ThirdHomomorphism}. The operator $\odot$ need not be
commutative, because word and paragraph order can matter. C-TreePO is
an audited learned approximation to this state-level discipline: the
summaries are learned states, and the local laws test whether those
states still support the task readout.

\paragraph{Lean artifact.}
The Lean artifact mirrors exactly this split. It proves the reusable
state-level algebra---build, merge, validity, tree closure, and final
readout---and separately exposes the large external facts as typed
theorem obligations. For HLL, the checked part is the
hash-to-register pipeline, pointwise-max merge, state-level collapse,
the Poisson-mixture series package, and the fixed-cardinality
consequence obtained from a depoissonization-transfer obligation; the
Mellin-transform estimates and analytic depoissonization proof behind
the estimator constants remain named obligations, with the Big-O
consequence of the relative standard-error asymptotic checked in Lean.
For Agarwal-style summaries, the state-level interface, Count-Min
contrast, Misra--Gries/GK surfaces, interval and range-space targets,
finite VC traces, randomized root-validity/readout nesting, and
hybrid-promotion bookkeeping are typed; the compact randomized discrepancy
construction remains an explicit citation obligation rather than a hidden
assumption, and $\varepsilon$-kernel material is adjacent background unless
separately cited. The standalone Lean entry point for this literature layer is
\texttt{FormalProbability.ML.MergeableSummaries.Literature}; C-TreePO
imports that surface and re-exports only the bridge names used by this
paper. The corresponding Lean names are:
\begin{center}
\small
\path|local_laws_of_sketch|\\
\path|mergeableSketchSummaryClass_subset_exactLocalLawSubspace|\\
\path|RelationalMergeablePreferenceShape|\\
\path|RandomizedRelationalMergeablePreferenceShape|\\
\path|CTreePOToAgarwalTransform|
\end{center}

\paragraph{The formal correspondence.}
The correspondence with C-TreePO is a qualified equivalence, not a bare
claim that every scalar answer has a merge rule. There are two cases.
The strict case is an oracle homomorphism, where the task value itself
has a merge operation. This is clean but too strong for many sketches:
scalar distinct count cannot merge $|A|$ and $|B|$ into $|A\cup B|$
without overlap information. The state-level case is the one used by
classical mergeable summaries. The state carries the missing
information, the state is merged locally, and the readout happens after
the root state has been formed.

\begin{proposition}[Strict and State-Level Mergeable Limits]\label{prop:mergeable-reduction-background}
Under Assumption~\ref{ass:context} (context-compatible oracle), suppose $g$ is deterministic and satisfies global versions of the preservation laws:
\[
\text{(A1)}\ \forall z,\ g(z) \sim z,\qquad
\text{(A2)}\ \forall u,v,\ u \concat v \sim g(g(u)\concat g(v)),
\]
Define the induced summary merge $\oplus_g(s,t):=g(s\concat t)$. Then:
\begin{enumerate}[nosep]
\item \emph{Strict oracle homomorphism.} If there exists
$m_{\Y}:\Y\times\Y\to\Y$ such that
$m_{\Y}(\fstar(u),\fstar(v))=\fstar(u\concat v)$ for all $u,v$, then
$\fstar$ is a homomorphic readout of the induced state: merging summaries and
then reading out gives the same oracle value as reading the raw concatenation,
and $\oplus_g$ is associative modulo $\sim$.
\item \emph{Classical state-level reduction.} More generally, let
$(\mathrm{encode},\oplus_B,\mathrm{query})$ be a bounded mergeable-summary
scheme for a query family $Q_q$ in the sense of
\citet{Agarwal2013MergeableTODS}. If $g$ serializes the sketch state
($g(x)=\mathrm{encode}(x)$ on raw shards and
$g(s\concat t)=s\oplus_B t$ on serialized states) and the C-TreePO readout is
$f(s,q)=\mathrm{query}(s,q)$, then the C-TreePO tree realizes the same
state-level mergeable summary under the induced merge $\oplus_g$.
\end{enumerate}
\end{proposition}

\noindent\textit{Proof intuition.} The strict case chains A1 and A2
through context compatibility to recover homomorphic readout and
associativity modulo the oracle. The classical case is state-level:
bounded sketch states compose by $\oplus_B$, and the readout is
applied only after the state has been merged. In the randomized
Agarwal-style case, the same state-level reduction is interpreted on the
root-validity event: with the paper's success probability, the merged root
state is valid for the union stream, and then the readout is correct. The
full proof is in
\S\ref{app:proof-mergeable} below.

\paragraph{Exact instantiation checklist.}
To use this proposition for a concrete mergeable summary, specify
(i) a state space $\mathcal{S}$, (ii) a validity relation $V(D,s)$,
(iii) a builder $\mathrm{build}(D)$, (iv) a state merge
$\mathrm{merge}(s_L,s_R)$, and (v) a readout
$\mathrm{readout}(s)$.  The required facts are
\[
V(D,\mathrm{build}(D)),\qquad
V(D_L,s_L)\wedge V(D_R,s_R)
\Rightarrow
V(D_L\concat D_R,\mathrm{merge}(s_L,s_R)),
\]
and
\[
V(D,s)\Rightarrow \mathrm{readout}(s)=F(D).
\]
C-TreePO implements these facts by choosing $g$ to build states at leaves and
merge serialized child states at internal nodes, and choosing $f$ to be the
root readout. C1 checks the leaf build case, C3 checks the merge case, and C2
is the additional re-use/stability condition needed when an already-produced
summary can itself be summarized again. None of these conditions says that
scalar child answers merge; scalar answer mergeability is a stricter special
case.

\paragraph{Variants the framework covers.}
Appendix~\ref{sec:hll-parity} exhibits this collapse numerically on one
representative mergeable sketch, HyperLogLog. Running a classical HLL
(native plus the Apache DataSketches reference) through the unified
\texttt{fit()} framework with $n_{\mathrm{leaves}} \in \{1, 2, 4, 8,
16\}$ yields byte-identical register states between TreePO-reduced and
flat pipelines for the register-based implementation, and
estimate-equivalence within the HLL relative standard error for the
DataSketches reference. The broader \texttt{treepo-bench suite
classical-sketches} path repeats the same tree-reduction comparison
for official DataSketches distinct-counting, frequency, quantile, and
set-operation sketches plus Tuple/VarOpt sampling checks, reporting
memory/error curves over capacities and leaf counts, merge-schedule
spread, bound coverage, and distance to the best official sketch
floor. It is executed through the same \texttt{fit()} interface as the
HLL parity experiment. This matters for the learned comparisons: the
scalar learned-$g{+}f$ companion can take an exact target or an
official sketch query function as $\fstar$, then learn both the merge
state and readout without forcing the hidden state to be byte-
compatible with the reference library. The same pipeline can swap its
oracle $\fstar$ from analytic cardinality to the HLL's own scoring
head, and under that oracle a fully end-to-end learned $g$ and $f$
track the classical sketch within a small multiple of the RSE floor;
a constrained variant in which $g$ must produce outputs readable by
the classical HLL estimator degrades with $L$, quantifying the cost of
the register-space constraint on depth. In the projection language,
the classical rows start inside the exact local-law subspace, while
the learned rows test different representable subspaces and readouts
for how closely training can reach that target. When $\fstar$ is an
official sketch query function, C-TreePO is being asked to recover
that classical sketch; when $\fstar$ is semantic, the same laws define
the learned generalization. The formal claim is not HLL-specific: any
sketch that supplies the same \texttt{SketchLeafPreserving},
\texttt{SketchMergeCompatible}, and \texttt{SketchSummaryCompatible}
witnesses enters the exact local-law subspace through
\texttt{mergeableSketchSummaryClass\_subset\_exactLocalLawSubspace}.

\subsection{Sketch Sufficiency Boundary}\label{app:proof-m-lt-k}

The following observation motivates the general mechanism: when a bounded sketch cannot represent the interaction the oracle depends on, the failure is structural and more data cannot fix it.

\begin{proposition}[Failure Boundary for $m < k$]\label{prop:m_lt_k}
For a threshold target $\tau_k(r) := \One\{C(r) \ge k\}$ on binary indicators $r \in \{0,1\}^n$ with $C(r) = \sum_j r_j$, if a top-$m$ sketch retains only $m < k$ indicators, there exist inputs with identical sketch state but different target values. No estimator restricted to that sketch class can recover $\tau_k$ without ambiguity.
\end{proposition}

\begin{proof}
Fix integers $k,m$ with $1 \le m < k$.
Construct two documents:
\[
r^{-} = (\underbrace{1,\dots,1}_{m\ \text{times}}, 0, \dots, 0), \qquad
r^{+} = (\underbrace{1,\dots,1}_{k\ \text{times}}, 0, \dots, 0),
\]
with common length $n \ge k$. Then $C(r^-) = m < k$ and $C(r^+) = k \ge k$, so
$\tau_k(r^-) = 0$ and $\tau_k(r^+) = 1$.
However, both sequences have at least $m$ unit entries, and all top-$m$ retained entries are ones.
Hence $S_m(r^-) = S_m(r^+)$.
Therefore two inputs with identical sketch state yield different target values.
\end{proof}

\subsection{Proof of Theorem~\ref{thm:one-pass} (Preservation)}\label{app:proof-preservation}

\begin{proof}
Fix a realized full binary tree $T$ over document $x$.
For each node $u$, define latent raw span $S(u)$ recursively by setting $S(u) = b$ when $u$ is a leaf containing block $b$, and setting $S(u) = S(u_L) \concat S(u_R)$ when $u$ is an internal node with children $(u_L, u_R)$.
Define stored summary $s_u$ recursively by setting $s_u = g(S(u))$ at leaves and $s_u = g(s_{u_L} \concat s_{u_R})$ at internal nodes.
Let $r$ be the root and set $Z^{(1)}(x) := s_r$.
Assume C1 and C3 hold almost surely on all realized calls.

We show $s_u \sim S(u)$ for every node $u$ by structural induction on node height.

\textbf{Base case.} If $u$ is a leaf, C1 gives $s_u = g(S(u)) \sim S(u)$.

\textbf{Inductive step.} Let $u$ be an internal node with children $u_L, u_R$ and assume $s_{u_L} \sim S(u_L)$ and $s_{u_R} \sim S(u_R)$.
By context compatibility (Assumption~\ref{ass:context}), applied first to right-concatenation with $s_{u_R}$ and then to left-concatenation with $S(u_L)$:
\[
s_{u_L} \concat s_{u_R} \sim S(u_L) \concat s_{u_R} \sim S(u_L) \concat S(u_R) = S(u).
\]
The remaining obligation is the realized merge-law check at node $u$. In the paper's two-link C3 notation, applying C3 to the child summaries gives the displayed equivalence below; together with the context-compatibility equivalence above, this is the internal-node zero-distortion conclusion. In Lean, L2 packages that composed conclusion directly. Thus
\[
g(s_{u_L} \concat s_{u_R}) \sim s_{u_L} \concat s_{u_R}.
\]
Since $s_u = g(s_{u_L} \concat s_{u_R})$, transitivity gives $s_u \sim S(u)$.

At root $r$, we have $s_r \sim S(r) = x$, hence
\[
\metric\!\big(\fstar(Z^{(1)}(x)), \fstar(x)\big) = 0
\]
almost surely, so $\E[\metric(\fstar(Z^{(1)}(x)), \fstar(x))] = 0$.

\textbf{Multi-round extension (add C2).}
Define rounds by $Z^{(R+1)}(x) := g(Z^{(R)}(x))$ for $R \ge 1$.
Because $Z^{(1)}(x) = s_r \in \operatorname{range}(g)$ and $g$ maps into its range, each $Z^{(R)}(x)$ is in $\operatorname{range}(g)$.
By C2: $Z^{(R+1)}(x) \sim Z^{(R)}(x)$ for all $R \ge 1$.
We prove $Z^{(R)}(x) \sim x$ by induction on $R$.
Base $R = 1$ holds from the nodewise equivalence argument above.
If $Z^{(R)}(x) \sim x$ and $Z^{(R+1)}(x) \sim Z^{(R)}(x)$, transitivity yields $Z^{(R+1)}(x) \sim x$.
Thus $\E[\metric(\fstar(Z^{(R)}(x)), \fstar(x))] = 0$ for all $R \ge 1$.
\end{proof}

\subsection{Proof of Proposition~\ref{prop:mergeable-reduction-background} (Strict and State-Level Mergeable Limits)}\label{app:proof-mergeable}

\begin{proof}
Assume Assumption~\ref{ass:context} and deterministic $g$ with A1--A2 as
stated in Proposition~\ref{prop:mergeable-reduction-background}. Define
\[
\oplus_g(s,t):=g(s\concat t).
\]
We first prove the strict oracle-homomorphism case, then the state-level
classical-sketch reduction.

\textbf{(i) Strict oracle homomorphism.}
Assume there is an oracle-level merge $m_{\Y}$ with
$m_{\Y}(\fstar(u),\fstar(v))=\fstar(u\concat v)$ for all $u,v$.
For any $u,v$,
\[
g(u\concat v)\sim \oplus_g(g(u),g(v))
  \iff g(u\concat v)\sim g(g(u)\concat g(v)).
\]
By A1 on $u\concat v$, $g(u\concat v)\sim u\concat v$.
By A2, $u\concat v\sim g(g(u)\concat g(v))$.
Transitivity yields merge-route equivalence. Because A1 gives
$g(u)\sim u$ and $g(v)\sim v$, the definition of $m_{\Y}$ gives
\[
m_{\Y}\big(\fstar(g(u)),\fstar(g(v))\big)
  = m_{\Y}\big(\fstar(u),\fstar(v)\big)
  = \fstar(u\concat v)
  = \fstar\big(\oplus_g(g(u),g(v))\big),
\]
where the last equality uses merge-route equivalence. Thus $\fstar$ is a
homomorphic readout of the induced state modulo $\sim$.

For any $a,b,c$, we need
\[
\oplus_g(\oplus_g(a,b),c)\sim \oplus_g(a,\oplus_g(b,c)),
\]
i.e.
\[
g(g(a\concat b)\concat c)\sim g(a\concat g(b\concat c)).
\]
By A1, outer $g$ can be removed on both sides modulo $\sim$, so it suffices to compare
$g(a\concat b)\concat c$ and $a\concat g(b\concat c)$.
Again by A1, $g(a\concat b)\sim a\concat b$ and $g(b\concat c)\sim b\concat c$.
Using context compatibility (Assumption~\ref{ass:context}), these imply
\[
g(a\concat b)\concat c \sim (a\concat b)\concat c,\qquad
a\concat g(b\concat c)\sim a\concat(b\concat c).
\]
Since concatenation is associative, $(a\concat b)\concat c = a\concat(b\concat c)$, hence both sides are equivalent to the same middle string, so they are equivalent to each other.
Therefore $\oplus_g$ is associative modulo oracle.

\textbf{(ii) Classical state-level reduction.}
Let $(\mathrm{encode},\oplus_B,\mathrm{query})$ be a bounded mergeable-summary
scheme for query family $Q_q$ in the sense of
\citet{Agarwal2013MergeableTODS}. Instantiate C-TreePO so that serialized
sketch states are valid summaries: on raw shards $g(x)=\mathrm{encode}(x)$,
on serialized states $g(s\concat t)=s\oplus_B t$, and
$f(s,q)=\mathrm{query}(s,q)$. Then for every pair of shards and every query,
\[
f(g(u\concat v),q)
  = \mathrm{query}(\mathrm{encode}(u\concat v),q)
  \approx
    \mathrm{query}(\mathrm{encode}(u)\oplus_B\mathrm{encode}(v),q)
  = f(\oplus_g(g(u),g(v)),q).
\]
The state remains bounded because the classical scheme is bounded after
merge. Associativity and schedule-invariance are inherited from the classical
merge operation, in the exact-state sense when $\oplus_B$ is exact and in the
readout-equivalence sense when the sketch guarantee is approximate. Therefore
the C-TreePO instantiation realizes the same state-level mergeable summary.
No oracle-level merge on $\Y$ is required; the readout can depend on state
that is not present in the scalar task value.

For randomized mergeable summaries, let $\Omega$ encode all random choices used
across the whole tree. Agarwal's randomized condition supplies a probability
lower bound for the event that the root state obtained by merging leaf states
is valid for the union stream. The C-TreePO readout is deterministic conditional
on this validity relation, so the readout-correctness event contains the
root-validity event and inherits the same probability lower bound. This is a
state-level nesting theorem, not a law for merging scalar query answers.
\end{proof}

\subsection{Kovachki Finite-Dimensionalization Lemmas}\label{app:kovachki-finite-dimensionalization}

The neural-operator step used in Appendix~\ref{sec:fno-primer} is not a
black box from ``continuous operator'' directly to ``certified
C-TreePO summarizer.'' It has three parts. Lemma~21 supplies compactness
for the union of finite-rank approximating images. Lemma~22 uses that
compactness to finite-dimensionalize the target operator on the compact
set. Theorems~11 and~13 then supply the neural-architecture
approximation statements. We record the C-TreePO-used Lemma~21/22 layer
below and route its output into the local-law bridge of
Appendix~\ref{app:approx-via-neural-operators}.

\paragraph{Kovachki Lemma 21, compact union.}
Let $X,Y$ be Banach spaces, $F:X\to Y$ continuous, $K\subset X$
compact, and $F_n:X\to Y$ continuous with
\[
\lim_{n\to\infty}\sup_{x\in K}\|F(x)-F_n(x)\|_Y=0.
\]
Then
\[
W=\bigcup_{n\ge 1}F_n(K)\cup F(K)
\]
is compact in $Y$. In the paper proof this is the device that lets one
apply uniform continuity of $G$ on one compact set containing both the
original compact $K$ and the finite-rank approximating images
$U_n^X(K)$. For the Lipschitz or globally uniformly continuous targets
we actually use in the C-TreePO bridge, the compact-union premise is
discharged directly by the uniform continuity of $G$ on the union $W$
above; the realized-call proof then proceeds without needing the
Lemma~21 stability hypothesis as a separate assumption.

\paragraph{Kovachki Lemma 22, finite-dimensionalization.}
Let $X,Y$ be Banach spaces with the approximation property and let
$G:X\to Y$ be continuous. For every compact $K\subset X$ and
$\varepsilon>0$, there are $J,J'\in\mathbb{N}$, continuous linear maps
$F_J:X\to\mathbb{R}^{J}$ and $G_{J'}:\mathbb{R}^{J'}\to Y$, and a
continuous map $\phi:\mathbb{R}^{J}\to\mathbb{R}^{J'}$ such that
\[
\sup_{x\in K}
\left\|G(x)-\bigl(G_{J'}\circ\phi\circ F_J\bigr)(x)\right\|_Y
\le \varepsilon.
\]
Moreover, $F_J$ can be written as a finite list of input functionals,
\[
F_J(x)=(w_1(x),\ldots,w_J(x)),\qquad w_j\in X^\ast,
\]
and $G_{J'}$ can be written as a finite expansion,
\[
G_{J'}(v)=\sum_{j=1}^{J'}v_j\beta_j,\qquad \beta_j\in Y.
\]
If $Y$ has a basis, the output representers can be chosen from a basis
extension. The two displayed forms (the finite-coordinate encoder $F_J$
and the finite expansion $G_{J'}$) are the finite-coordinate
presentation of the approximation property; this is equivalent to the
standard finite-rank-operator definition of AP. When the target $G$ is
uniformly continuous, the Lemma~21 stability argument discharges
automatically; the constructor then produces the full
$G_{J'}\circ\phi\circ F_J$ witness with its uniform compact error bound
as a direct corollary.

\paragraph{How this enters C-TreePO.}
For each realized call site, Lemma~22 produces a finite-dimensionalized
operator with a chosen tolerance, which is exactly the input the
uniform-on-compact approximation interface consumes. In the common
uniformly continuous case, AP, compactness, and uniform continuity
together yield that interface directly. The same witness is recorded
separately at each realized leaf, merge, and on-range call, after the
Banach-space representation has been transported to that call's domain
and codomain. Composing these witnesses with
Proposition~\ref{prop:neural-operator-bridge} gives approximate
theorem-backedness. Thus Lemma~22 helps the proof only up to the point
of producing realized-call approximation budgets; C1--C3 and the
audit/local-law theorems still perform the certification step.

\subsection{Neural-Operator Realization-to-Local-Law Bridge}\label{app:approx-via-neural-operators}

Appendix~\ref{sec:fno-primer} isolates the only part of the neural-operator
literature this paper uses. The external approximation theorem supplies an
ideal-to-realized approximation witness on compact realized-call sets. The
finite-dimensionalization part of that route is the Lemma~21/22 layer
spelled out in Appendix~\ref{app:kovachki-finite-dimensionalization}; the
equation-(6) architecture approximation remains the named external
Theorem~11/13 input.

Let $s^\dagger$ be an ideal deterministic summarizer on a fixed tree $T$,
and suppose $s^\dagger$ is exact theorem-backed. Choose three compact
realized-call sets: one for realized leaves, one for realized merge calls,
and one for the root-level on-range re-summary call. A Section~9-style
neural-operator approximation theorem is then invoked only to assert the
existence of a realized operator $s$ with uniform error at most
$\varepsilon$ on those compact sets.

What does \emph{not} happen is an automatic jump from operator error to
the paper's approximate local laws. The audit layer is phrased in terms
of violation probabilities, so one still needs explicit transfer
assumptions modelling how a uniform approximation tolerance $\varepsilon$
propagates to leaf, merge, and idempotence budgets
$(\varepsilon_{\mathrm{leaf}}, \varepsilon_{\mathrm{merge}},
\varepsilon_{\mathrm{idemp}})$. Proposition~\ref{prop:neural-operator-bridge}
packages this propagation: the approximate-local-law bundle it produces
is exactly the three displayed budgets plus a proof that the realized
operator satisfies C1/C2/C3 up to those budgets.

Once the bridge has produced an approximate-local-law bundle, the rest
of the proof route is exactly the one already used elsewhere in the
paper. The approximate-gap theorem gives
\[
\Delta_R \;\le\;
\varepsilon_{\mathrm{leaf}}
\;+\;
\varepsilon_{\mathrm{merge}}
\;+\;
(R-1)\,\varepsilon_{\mathrm{idemp}},
\]
and Theorem~\ref{thm:unified-gap} then transports this distortion budget to
the method-specific objective gap. So the role of the neural-operator theory
is precise: it is the realizability argument for obtaining a summarized
operator with a chosen upstream tolerance on the realized call compacts.
The local-law and audit machinery still does the certification work.
The regularized training surface uses the same bridge through
\texttt{oracleProjectionObjective\_le\_of\_uniformApproxExactTheoremBacked},
which is why the objective in Equation~\eqref{eq:oracle-projection-objective}
can be read as oracle approximation plus projection toward the local-law
subspace rather than as a new proof route. Appendix~\ref{app:projection-interpretation}
records the accompanying iff: this projection reading is equivalent to
saying that the approximation error is structured so that its zero set is
exactly the exact local-law subspace.

\subsection{Supporting Lemmas for the Gap Bound}\label{app:supporting-lemmas}

Before proving the main equivalence and gap theorems, we separate the two quantitative roles. Appendix~\ref{app:approx-via-neural-operators} is the upstream representation step: external neural-operator approximation plus transfer moduli produces a local-law distortion budget. The lemmas below are the downstream transport step: they show how a preference objective converts that oracle distortion budget into objective units.

\begin{lemma}[Sigmoid is 1-Lipschitz]\label{lem:sigmoid-lip}
$|\sigma(t_1) - \sigma(t_2)| \le |t_1 - t_2|$ for all $t_1, t_2 \in \R$, where $\sigma(t) = (1 + e^{-t})^{-1}$. Standard calculus: $|\sigma'(t)| = \sigma(t)(1-\sigma(t)) \le 1/4$.
\end{lemma}

\begin{lemma}[$-\log\sigma$ is 1-Lipschitz]\label{lem:neglogsig-lip}
$|{-}\log\sigma(t_1) - ({-}\log\sigma(t_2))| \le |t_1 - t_2|$ for all $t_1, t_2 \in \R$. Standard calculus: $|\frac{d}{dt}(-\log\sigma(t))| = |1 - \sigma(t)| < 1$.
\end{lemma}

\begin{lemma}[DPO pointwise loss is Lipschitz]\label{lem:dpo-lip}
Define the DPO pointwise loss for a preferred--dispreferred pair $(a_w, a_l)$:
\[
\ell_{\mathrm{DPO}}(x, a_w, a_l) := -\log\sigma\!\left(\beta\left[\log\frac{\pi_\theta(a_w|x)}{\pi_{\mathrm{ref}}(a_w|x)} - \log\frac{\pi_\theta(a_l|x)}{\pi_{\mathrm{ref}}(a_l|x)}\right]\right).
\]
If the policy log-ratio $x \mapsto \log\frac{\pi_\theta(a|x)}{\pi_{\mathrm{ref}}(a|x)}$ is $L_{\mathrm{pol}}$-Lipschitz in $\metric(\fstar(x), \fstar(x'))$ for every action $a$, then
\[
|\ell_{\mathrm{DPO}}(x, a_w, a_l) - \ell_{\mathrm{DPO}}(x', a_w, a_l)| \le 2|\beta|\, L_{\mathrm{pol}} \cdot \metric(\fstar(x), \fstar(x')).
\]
\end{lemma}

\begin{proof}
Define the logit difference $\Delta(x) := \log\frac{\pi_\theta(a_w|x)}{\pi_{\mathrm{ref}}(a_w|x)} - \log\frac{\pi_\theta(a_l|x)}{\pi_{\mathrm{ref}}(a_l|x)}$. By the triangle inequality on the two log-ratio terms, each $L_{\mathrm{pol}}$-Lipschitz:
\[
|\Delta(x) - \Delta(x')| \le 2 L_{\mathrm{pol}} \cdot \metric(\fstar(x), \fstar(x')).
\]
Composing: $\ell_{\mathrm{DPO}}(x, a_w, a_l) = (-\log\sigma)(\beta \cdot \Delta(x))$. By Lemma~\ref{lem:neglogsig-lip}, $-\log\sigma$ is 1-Lipschitz, and scaling by $|\beta|$ gives the claimed constant $2|\beta| L_{\mathrm{pol}}$.
\end{proof}

\begin{lemma}[DPO loss is oracle-measurable]\label{lem:dpo-oracle-meas}
If $\pi_\theta(a|\cdot)$ and $\pi_{\mathrm{ref}}(a|\cdot)$ are oracle-measurable---meaning $\metric(\fstar(x), \fstar(x')) = 0$ implies $\pi_\theta(a|x) = \pi_\theta(a|x')$ and $\pi_{\mathrm{ref}}(a|x) = \pi_{\mathrm{ref}}(a|x')$ for all $a$---and the pair generator is oracle-indexed---meaning $\metric(\fstar(x), \fstar(x')) = 0$ implies $\mathrm{gen}(x) = \mathrm{gen}(x')$---then the expected DPO loss $\E_{(a_w,a_l) \sim \mathrm{gen}(x)}[\ell_{\mathrm{DPO}}(x, a_w, a_l)]$ depends on $x$ only through $\fstar(x)$.
\end{lemma}

\begin{proof}
If $\metric(\fstar(x), \fstar(x')) = 0$, then by oracle-measurability of the policies, $\pi_\theta(a|x) = \pi_\theta(a|x')$ and $\pi_{\mathrm{ref}}(a|x) = \pi_{\mathrm{ref}}(a|x')$ for all $a$, so $\ell_{\mathrm{DPO}}(x, a_w, a_l) = \ell_{\mathrm{DPO}}(x', a_w, a_l)$ for all pairs. Since $\mathrm{gen}(x) = \mathrm{gen}(x')$, the expectation over pairs is identical.
\end{proof}

\begin{lemma}[Zero distortion implies pointwise equality on support]\label{lem:zero-dist-support}
Let $\mu_Z$ be the distribution of $Z^{(R)}(x)$ for fixed $x$, and suppose $\E[\metric(\fstar(Z^{(R)}(x)), \fstar(x))] = 0$. Since $\metric \ge 0$, we have $\metric(\fstar(z), \fstar(x)) = 0$ for all $z$ in the support of $\mu_Z$.
\end{lemma}

\begin{proof}
Suppose for contradiction that $\metric(\fstar(z_0), \fstar(x)) > 0$ for some $z_0$ with $\mu_Z(z_0) > 0$. Then
\[
\E[\metric(\fstar(Z), \fstar(x))] = \sum_z \mu_Z(z)\,\metric(\fstar(z), \fstar(x)) \ge \mu_Z(z_0) \cdot \metric(\fstar(z_0), \fstar(x)) > 0,
\]
contradicting $\E[\metric(\fstar(Z), \fstar(x))] = 0$.
\end{proof}

\subsection{Proof of Theorem~\ref{thm:pref-equiv} and Theorem~\ref{thm:unified-gap} (Equivalence and Gap Bounds)}\label{app:proof-equivalence}

We prove the results in increasing generality: first the method-agnostic exact equivalence, then the unified gap bound, then the DPO transport instantiation. When a neural-operator realization theorem is available, its output enters this proof only through the already-formed $\Delta_R$ budget.

\paragraph{Step 1: Method-agnostic exact equivalence.}

\begin{proof}[Proof of Theorem~\ref{thm:pref-equiv}]
Fix a preference learning method with population integrand $\mathcal{J}_\theta(x) = \E_{a \sim \mathrm{gen}(x)}[\ell(x, a)]$, where $\ell$ is the method-specific loss and $\mathrm{gen}$ is the pair or group generator.

Assume (i) oracle-measurability of $\ell$: $\metric(\fstar(x), \fstar(x')) = 0$ implies $\ell(x,a) = \ell(x',a)$ for all $a$; and (ii) oracle-indexing of the generator: $\metric(\fstar(x), \fstar(x')) = 0$ implies $\mathrm{gen}(x) = \mathrm{gen}(x')$. Together these give oracle-measurability of the integrand: there exists $\phi_\theta: \Y \to \R$ such that $\mathcal{J}_\theta(x) = \phi_\theta(\fstar(x))$.

Assume also Assumption~\ref{ass:context} and that C1--C3 hold on realized calls. By Theorem~\ref{thm:multi-round}, $\E[\metric(\fstar(Z^{(R)}), \fstar(x))] = 0$. By Lemma~\ref{lem:zero-dist-support}, $\metric(\fstar(z), \fstar(x)) = 0$ for all $z$ in the support of $Z^{(R)}$. Therefore
\[
\mathcal{J}_\theta(z) = \phi_\theta(\fstar(z)) = \phi_\theta(\fstar(x)) = \mathcal{J}_\theta(x)
\]
for all $z$ in the support. Taking expectations:
\[
\E[\mathcal{J}_\theta(Z^{(R)})] = \E[\mathcal{J}_\theta(X)].
\]
Since this holds for all $\theta$, the argmin sets are identical.

For DPO, oracle-measurability of the loss follows from Lemma~\ref{lem:dpo-oracle-meas}. For GRPO-PL, the argument is identical with the Plackett-Luce ranking loss $\ell_{\mathrm{PL}}(x, \sigma) = -\sum_{i=1}^k \log \frac{\exp(s_{\sigma(i)}(x))}{\sum_{j \ge i} \exp(s_{\sigma(j)}(x))}$ in place of the Bradley-Terry loss, and with the ranker assumed oracle-indexed. For GRPO-RL, the objective includes clipped surrogate ratios and a KL penalty, but both depend on $x$ only through oracle-measurable policies and rewards, so the same factorization applies.
\end{proof}

\paragraph{Step 2: Unified gap bound.}

\begin{proof}[Proof of Theorem~\ref{thm:unified-gap}]
Let $\mu_X$ and $\mu_Z$ denote the distributions of $X$ and $Z^{(R)}$ respectively. Define the expected integrand $E_\theta(x) := \E_{a \sim \mathrm{gen}(x)}[\ell(x, a)]$.

Assume an $L$-Lipschitz condition: $|E_\theta(x) - E_\theta(z)| \le L \cdot \metric(\fstar(x), \fstar(z))$ for all $x, z$.

Because $Z = Z^{(R)}(X)$ is a function of $X$, we can write the gap using the conditional distribution $\mu_{Z|X}$:
\begin{align*}
\E_X[E_\theta(X)] - \E_Z[E_\theta(Z)]
&= \sum_x \mu_X(x) \left[ E_\theta(x) - \sum_z \mu_{Z|X=x}(z)\, E_\theta(z) \right].
\end{align*}
Taking absolute values and applying the Lipschitz bound pointwise:
\begin{align*}
\left|\E_X[E_\theta(X)] - \E_Z[E_\theta(Z)]\right|
&\le \sum_x \mu_X(x) \sum_z \mu_{Z|X=x}(z)\, |E_\theta(x) - E_\theta(z)| \\
&\le L \sum_x \mu_X(x) \sum_z \mu_{Z|X=x}(z)\, \metric(\fstar(x), \fstar(z)) \\
&= L \cdot \E\!\left[\metric(\fstar(X), \fstar(Z^{(R)}))\right] =: L \cdot \Delta_R.
\end{align*}
The interchange of sum and absolute value uses absolute convergence, which holds because $E_\theta$ is bounded.
\end{proof}

\noindent If the coupled distortion was produced by the neural-operator route, combine this with Proposition~\ref{prop:neural-operator-bridge}. Writing the transfer moduli as $\omega_{\mathrm{leaf}},\omega_{\mathrm{merge}},\omega_{\mathrm{idemp}}$, the bridge yields
\[
\Delta_R
\le
\omega_{\mathrm{leaf}}(\varepsilon)
+\omega_{\mathrm{merge}}(\varepsilon)
+(R-1)\omega_{\mathrm{idemp}}(\varepsilon),
\]
so the same proof gives
\[
|G_{\mathrm{meth}}|
\le
L\Big(\omega_{\mathrm{leaf}}(\varepsilon)
+\omega_{\mathrm{merge}}(\varepsilon)
+(R-1)\omega_{\mathrm{idemp}}(\varepsilon)\Big).
\]
The observational certificate of Theorem~\ref{thm:e2e} then replaces or supplements this design target with the sampled estimate and its calibration/sampling envelopes.

\paragraph{Step 3: DPO transport constant.}

Instantiating the unified gap bound for DPO requires verifying only the downstream transport condition on $E_\theta$; the summarizer's representational error has already been converted into $\Delta_R$ by the neural-operator/local-law or audit route. Lemma~\ref{lem:dpo-lip} establishes that the pointwise DPO loss is $2|\beta|L_{\mathrm{pol}}$-Lipschitz. When the generator is oracle-indexed (hence constant on oracle equivalence classes), the expected integrand inherits the same Lipschitz constant: for any $x, x'$ with $\mathrm{gen}(x) = \mathrm{gen}(x')$,
\[
|E_\theta(x) - E_\theta(x')| = \left|\sum_{(a_w, a_l)} \mathrm{gen}(x)(a_w, a_l)\, [\ell_{\mathrm{DPO}}(x, a_w, a_l) - \ell_{\mathrm{DPO}}(x', a_w, a_l)]\right| \le 2|\beta|\,L_{\mathrm{pol}} \cdot \metric(\fstar(x), \fstar(x')),
\]
where the last step uses Lemma~\ref{lem:dpo-lip} inside the sum and the fact that the generator weights sum to~1.

For GRPO-PL and GRPO-RL, the Lipschitz condition on the expected group loss enters as the explicit assumption \texttt{ExpectedGroupLossLipschitz}: the expected loss over group samples is Lipschitz in oracle distance, with the constant determined by the method-specific structure (Plackett-Luce ranking for GRPO-PL; clipped surrogate plus KL penalty for GRPO-RL). It is a hypothesis passed to the gap theorems, and the formalization proves sufficient conditions for it from first principles: a finite fixed-ranker Plackett--Luce condition for GRPO-PL and a finite pointwise Lipschitz condition for GRPO-RL (\texttt{OPT/RUMSufficientConditions.lean}). The Random Utility Model (Appendix~\ref{sec:verification}) is one way to motivate the assumption when those sufficient conditions are not directly checkable.

\subsection{Proof of Theorem~\ref{thm:e2e} (Audited Preference-Gap Bound)}\label{app:proof-e2e}

The audited preference-gap bound composes the representation/local-law budget with two independently established statistical steps: design-based unbiasedness of the distortion estimator, and deterministic method transport from distortion to the preference gap. Analytic neural-operator bounds can set a design target for the budget; the audit estimates the budget actually realized by the deployed tree.

\begin{proof}
\textbf{Part (a): Distortion objects and HT unbiasedness.}
Fix a realized tree-population $\mathcal{U}=\bigcup_x U(x)$ with $N := |\mathcal{U}|$, where $U(x)$ contains \emph{all realized nodes} of document $x$ (including the root). For each unit $i \in \mathcal{U}$, define
\[
d_i^{\star}:=\metric\!\big(\fstar(s_i),\fstar(S(i))\big),\qquad
d_i^{J}:=\metric\!\big(f(s_i),f(S(i))\big),
\]
and corresponding finite-population means
\[
\mu_{\mathrm{dist}}^{\star}:=\frac{1}{N}\sum_{i\in\mathcal{U}} d_i^{\star},\qquad
\mu_{\mathrm{dist}}^{J}:=\frac{1}{N}\sum_{i\in\mathcal{U}} d_i^{J}.
\]
Let $Z_i \in \{0,1\}$ be the inclusion indicator with propensity $\pi_i := \Prob(Z_i=1\mid\mathcal{F})$ and $0<\pi_i\le 1$. The unclipped HT estimator for judge distortion is
\[
\hat{\mu}_{\mathrm{HT}}^{J}:=\frac{1}{N}\sum_{i\in\mathcal{U}} \frac{Z_i}{\pi_i}d_i^{J}.
\]
Conditional unbiasedness is immediate:
\[
\E\!\left[\hat{\mu}_{\mathrm{HT}}^{J}\mid\mathcal{F}\right]
=\frac{1}{N}\sum_{i\in\mathcal{U}}\frac{\E[Z_i\mid\mathcal{F}]}{\pi_i}d_i^{J}
=\frac{1}{N}\sum_{i\in\mathcal{U}} d_i^{J}
=\mu_{\mathrm{dist}}^{J}.
\]

The full-document label $y_{\mathrm{doc}}=f^\ast(x)$ is a separate complete-label object and is not itself the node-population estimand unless the top node is explicitly included in $\mathcal{U}$ through the node-sampling law. Likewise the final-summary/tree-side quantity $y_{\mathrm{tree}}$ is a different object from $y_{\mathrm{doc}}$ except under exact preservation or exact-control assumptions. The HT statement here concerns only the sampled-node finite-population mean over $\mathcal{U}$; the document-level top loss is evaluated directly from the complete document labels.

\textbf{Part (b): Method transport.}
By Theorem~\ref{thm:unified-gap}, for the method-specific constant $C_{\mathrm{meth}}$,
\[
|G_{\mathrm{meth}}| \le C_{\mathrm{meth}}\,\mu_{\mathrm{dist}}^{\star}.
\]

\textbf{Part (c): Envelope decomposition.}
Let $\hat{\mu}_{\mathrm{dist}}$ denote the deployed distortion estimator (possibly clipped). Suppose on an event $\mathcal{E}$ we have the three envelope inequalities:
\[
C_{\mathrm{meth}}\!\left|\mu_{\mathrm{dist}}^{\star}-\mu_{\mathrm{dist}}^{J}\right|\le B_{\mathrm{cal}},\quad
C_{\mathrm{meth}}\!\left|\mu_{\mathrm{dist}}^{J}-\hat{\mu}_{\mathrm{HT}}^{J}\right|\le B_{\mathrm{est}},\quad
C_{\mathrm{meth}}\!\left|\hat{\mu}_{\mathrm{HT}}^{J}-\hat{\mu}_{\mathrm{dist}}\right|\le B_{\mathrm{clip}}.
\]
Then on $\mathcal{E}$,
\begin{align*}
|G_{\mathrm{meth}}|
&\le C_{\mathrm{meth}}\,\mu_{\mathrm{dist}}^{\star} \\
&\le C_{\mathrm{meth}}\,\hat{\mu}_{\mathrm{dist}}
 + C_{\mathrm{meth}}\!\left|\mu_{\mathrm{dist}}^{\star}-\mu_{\mathrm{dist}}^{J}\right|
 + C_{\mathrm{meth}}\!\left|\mu_{\mathrm{dist}}^{J}-\hat{\mu}_{\mathrm{HT}}^{J}\right|
 + C_{\mathrm{meth}}\!\left|\hat{\mu}_{\mathrm{HT}}^{J}-\hat{\mu}_{\mathrm{dist}}\right| \\
&\le C_{\mathrm{meth}}\,\hat{\mu}_{\mathrm{dist}} + B_{\mathrm{cal}} + B_{\mathrm{est}} + B_{\mathrm{clip}}.
\end{align*}
The middle step is the telescoping decomposition
\[
\mu_{\mathrm{dist}}^{\star}
=\hat{\mu}_{\mathrm{dist}}
+(\mu_{\mathrm{dist}}^{\star}-\mu_{\mathrm{dist}}^{J})
+(\mu_{\mathrm{dist}}^{J}-\hat{\mu}_{\mathrm{HT}}^{J})
+(\hat{\mu}_{\mathrm{HT}}^{J}-\hat{\mu}_{\mathrm{dist}})
\]
followed by triangle inequality and $\hat{\mu}_{\mathrm{dist}}\ge 0$ (nonnegative distortions with nonnegative weights). This is exactly the stated audited bound.

\textbf{Part (d): High-probability form.}
If $\Prob(\mathcal{E}) \ge 1-\delta$, the same inequality holds with probability at least $1-\delta$. In practice, $\mathcal{E}$ is formed by intersecting concentration/calibration events; e.g., when separate events $\mathcal{E}_{\mathrm{cal}},\mathcal{E}_{\mathrm{est}},\mathcal{E}_{\mathrm{clip}}$ are used, one may set $\mathcal{E}:=\mathcal{E}_{\mathrm{cal}}\cap\mathcal{E}_{\mathrm{est}}\cap\mathcal{E}_{\mathrm{clip}}$ and apply a union bound to obtain an explicit confidence level.
\end{proof}

\noindent\textit{Formalization.} The Lean artifact proves this composition end to end over the independent Bernoulli sampling design. The method certificates (\texttt{dpo\_treepo\_end\_to\_end\_certificate} and the GRPO variants) supply parts~(a)--(b), and \texttt{dpo\_treepo\_certificate\_instantiates\_error\_stack} instantiates the abstract certificate stack (\texttt{PaperErrorStack}) with that concrete probability space: the estimation and clipping deviations are the realized HT estimator minus its mean, with explicit Hoeffding failure probabilities discharging the corresponding events. The realized display of parts~(c)--(d) is \texttt{dpo\_treepo\_realized\_estimator\_certificate}: for any slack $t>0$, with probability at least $1-2\exp\!\big(-t^2/(8N(D_{\max}/\pi_{\min})^2)\big)$ under the sampling design, $|G_{\mathrm{meth}}| \le C_{\mathrm{meth}}(\hat{\mu}_{\mathrm{HT}} + t)$, where $\hat{\mu}_{\mathrm{HT}}$ is the realized unclipped HT distortion estimate. In the formal model the audited distortion is evaluated by $\fstar$ itself, so the judge-calibration envelope enters as a hypothesis slot ($B_{\mathrm{cal}}>0$ certified with zero failure probability) and the concentration slack is carried in distortion units in the stack instance; the fully transported display is the realized-estimator theorem. GRPO-PL analogues are \texttt{grpo\_pl\_treepo\_realized\_estimator\_certificate} and \texttt{grpo\_pl\_treepo\_certificate\_instantiates\_error\_stack}, and GRPO-RL analogues (with the shared pointwise-Lipschitz constant $L$ of the clipped surrogate plus KL loss as the transport constant) are \texttt{grpo\_rl\_treepo\_realized\_estimator\_certificate} and \texttt{grpo\_rl\_treepo\_certificate\_instantiates\_error\_stack} (\texttt{DSL/TreePOEndToEndGlue.lean}).

\subsection{Audit Robustness: Adversarial Sampling}\label{app:audit-robustness}

The HT guarantee is robust to non-benign sampling designs: unbiasedness does not require friendly sampling, only correct logged marginal propensities and strict positivity. The Lean artifact formalizes this arbitrary-design statement directly and then separately records the variance side condition needed for constrained designs. This supports constrained worst-case analysis without treating the sampling law as benign.

Define a design class
\[
\mathcal{D}(\pi_{\min},\kappa)
:=\left\{\pi:\ \pi_i\in[\pi_{\min},\kappa\pi_{\min}],\ \sum_{i\in\mathcal{U}}\pi_i=n\right\},
\]
which constrains the antagonist by a minimum inclusion probability and a propensity spread cap. Within this class, an adversary can increase variance by skewing mass toward high-risk units, but cannot alter the target expectation when HT reweighting uses the true logged $\pi_i$.

For independent Bernoulli inclusion, or more generally for a design whose covariance structure is bounded by the same Bernoulli proxy, bounded distortions $0\le Y_i\le D_{\max}$ imply
\[
\mathrm{Var}\!\left(\hat{\mu}_{\mathrm{HT}}\mid\mathcal{F}\right)
\le
\frac{D_{\max}^2}{N}\cdot \frac{1-\pi_{\min}}{\pi_{\min}},
\]
equivalently $\frac{D_{\max}^2}{N}(1/\pi_{\min}-1)$, so $\pi_{\min}$ is the primary robustness lever. A clipping cap $w_{\max}$ gives the corresponding operational variant by replacing uncontrolled inverse-propensity weights with the reported cap and recording the induced clipping term in the certificate. Practically, certificates should therefore be reported over a stress-test grid of $(\pi_{\min}, w_{\max})$ values rather than for a single nominal design.
